\begin{abstract}
Salo proved that every unishift has dense totally periodic points by recoding it, after an integral change of basis and convex blocking, as an inductive uniform binondeterministic spacetime.  We make the periodic-completion step quantitative for a fixed two-dimensional presentation.  If the one-step horizontal trace uses alphabet $A$, both nondeterministic directions have horizontal radius at most $R$, and a globally valid $W\times H$ rectangle is prescribed, then it occurs in a doubly periodic point with horizontal and vertical periods
\[
 m\leq W+2R(H-1)+\card A-1,
 \qquad
 T\leq H+\card A^m.
\]
The proof closes the bottom dependency cone in the balanced trace graph and then closes the resulting path in the balanced finite graph of $m$-periodic rows.  Both closings are constructive.  We give an explicit exponential-time algorithm, a nontrivial finite check, and a presentation-local consequence: every nonidentity radius-$\rho$ automorphism is separated by its action on a computably enumerable rectangular periodic-point set whose periods satisfy explicit bounds.  Thus the usual residual-finiteness argument becomes quantitative.  The constants are not asserted in the original coordinates of an arbitrary unishift: the basis change, blocking alphabet, and presentation radius are genuine costs.
\end{abstract}
